\documentclass[10pt,twocolumn]{article}
\usepackage[T1]{fontenc}
\usepackage{lmodern}

\usepackage[margin=1in]{geometry}
\usepackage{amsmath,amssymb}
\usepackage{graphicx}
\usepackage{booktabs}
\usepackage{multirow}
\usepackage{hyperref}
\usepackage{microtype}
\usepackage{xcolor}
\usepackage{url}
\usepackage{cite}
\usepackage{caption}
\usepackage{subcaption}
\usepackage{float}

\hypersetup{
  colorlinks=true,
  linkcolor=blue,
  citecolor=blue,
  urlcolor=blue
}

\title{\textbf{Physiological Stress Detection from Wrist-Worn Sensors:\\
Comparing Random Forests, 1D CNNs, and Transformers\\
with Systematic Channel Ablation on WESAD}}

\author{
  Ranga Narayanaswami, PhD\\
  BeaconPoint AI Systems\\
  \texttt{ranga@beaconpointai.com}
}

\date{}

\begin{document}

\maketitle

\begin{abstract}
We evaluate three model families---Random Forest (RF) with handcrafted features,
a 1D Convolutional Neural Network (CNN), and a patch-based Transformer---for
binary stress detection using the WESAD dataset with
Leave-One-Subject-Out Cross-Validation (LOSO-CV).
All models share the same six-channel wrist-sensor input
(BVP, ACC$_{x/y/z}$, TEMP, EDA) derived from the Empatica E4.
Under strict LOSO-CV, RF achieves a mean F1 of $0.831 \pm 0.292$,
the CNN $0.897 \pm 0.218$, and the Transformer $0.893 \pm 0.211$;
the two deep models are statistically indistinguishable (Wilcoxon
signed-rank $p=1.00$) and both outperform RF, demonstrating that deep
architectures provide modest but consistent gains over handcrafted
features on this 15-subject dataset.
Beyond the model comparison, our primary contribution is a systematic
\emph{channel ablation study}: we isolate the contribution of each
wrist-sensor modality across both CNN and RF architectures.
Electrodermal activity (EDA) emerges as the dominant stress biomarker,
contributing a $+0.295$ F1 gain when added to a BVP$+$ACC baseline
($0.589 \to 0.884$) in the CNN.
Accelerometer (ACC) signals, by contrast, \emph{decrease} stress-discrimination
performance ($-0.148$ F1) in the sedentary WESAD Trier Social Stress Test (TSST)
protocol, where subjects are seated and motion is near-constant.
Skin temperature (TEMP) provides negligible marginal benefit as a direct
stress discriminator once EDA is included.
These findings carry practical implications for wearable sensor design:
a two-channel BVP$+$EDA configuration achieves $0.857 \pm 0.179$ F1,
within 2 points of the full six-channel model while eliminating two sensor
modalities in sedentary contexts.
Importantly, ACC and TEMP retain value as \emph{context signals} in
free-living deployment: ACC can detect vigorous exercise and gate stress
classification during physical activity (when thermoregulatory EDA
confounds the stress signal), while TEMP can flag elevated ambient conditions
that similarly elevate EDA through sweating rather than sympathetic arousal.
Per-subject analysis reveals clinically meaningful inter-individual
variability, including subjects for whom EDA reduces classification accuracy (inverted response)
or TEMP actively confounds classification.
\end{abstract}

\section{Introduction}
\label{sec:intro}

Psychological stress is a major determinant of cardiovascular disease,
metabolic syndrome, and mental health outcomes~\cite{yaribeygi2017}.
Continuous, unobtrusive stress monitoring via consumer wearable devices---
smartwatches and fitness bands---offers the possibility of population-scale
early detection and intervention. The Empatica E4, a research-grade wrist
device, captures four physiological modalities simultaneously:
blood volume pulse (BVP), measured via photoplethysmography (PPG), at 64\,Hz,
three-axis accelerometry (ACC) at 32\,Hz,
electrodermal activity (EDA) at 4\,Hz, and skin temperature (TEMP) at 4\,Hz.
These signals, and sensors similar to them, appear increasingly in consumer
devices: optical heart rate monitors are nearly universal, EDA sensors appear
in Fitbit Sense 2 and Samsung Galaxy Watch, and skin temperature sensors
appear in the Apple Watch Ultra and Fitbit Charge 6.

Yet despite substantial literature on physiological stress detection,
two questions remain underexplored in the wrist-sensor context:

\begin{enumerate}
  \item \textbf{Architecture}: Do deep learning models (CNNs, Transformers)
  substantially outperform traditional RF classifiers on multimodal wrist signals
  under realistic evaluation conditions?

  \item \textbf{Channel necessity}: Which sensor modalities actually drive
  performance? Can a subset of channels match the full-sensor baseline?
\end{enumerate}

Most published comparisons use either within-subject validation
(training and testing on windows from the same subjects), which inflates
reported accuracy, or are conducted on chest-worn clinical sensors
(ECG, EMG, respiration) that are unavailable in consumer devices~\cite{mekruksavanich2021,cross_modality_2025}.
We restrict our study to wrist signals and adopt strict LOSO-CV throughout,
holding out one subject entirely for testing at each fold.

Our contributions are:
\begin{itemize}
  \item A head-to-head comparison of RF, 1D CNN, and patch-based Transformer
  on the WESAD benchmark under LOSO-CV with identical train/test splits
  and identical six-channel input.
  \item A systematic CNN channel ablation study spanning six configurations
  from BVP-only (1 channel) to the full six-channel stack, quantifying
  the marginal value of ACC, TEMP, and EDA.
  \item A parallel RF feature-group ablation across ten configurations
  (HRV, ACC, EDA, TEMP, and combinations), confirming that EDA dominance
  is a property of the physiological signal rather than a modeling artifact.
  \item Per-subject analysis of outlier subjects revealing clinically meaningful
  inter-individual variability in physiological stress responses.
  \item Discussion of deployment implications: thermoregulatory EDA confounds,
  ACC-based activity gating, and device sensor selection for stress monitoring.
\end{itemize}

\section{Related Work}
\label{sec:related}

\subsection{WESAD Benchmark Studies}

Schmidt et al.~\cite{schmidt2018wesad} introduced WESAD, evaluating LDA,
RF, AdaBoost, and linear discriminant analysis on HRV and physiological
features, reporting macro-averaged F1 scores of approximately 0.93 using
both chest and wrist sensors with leave-one-subject-out evaluation.
Their analysis included wrist-only subsets, establishing the dataset as
the standard benchmark for wrist-sensor stress research.

Subsequent work has explored the WESAD benchmark extensively.
Mekruksavanich et al.~\cite{mekruksavanich2021} applied deep LSTM networks
to WESAD signals, reporting high accuracy but using train/validation splits
that allow windows from the same subject to appear in both sets---a protocol
that significantly inflates performance estimates relative to LOSO-CV.

A recent cross-modality study~\cite{cross_modality_2025} applied
transformer-based architectures to WESAD, reporting validation accuracies
above 99\%. However, this work uses \emph{chest}-mounted sensors
(ECG, EMG, respiration)---modalities unavailable in consumer wearables---
and evaluates with within-subject validation rather than LOSO-CV.
Their research question (cross-modal knowledge transfer between physiological
modalities) is orthogonal to ours; we cite it to underscore the performance
gap attributable to sensor placement and evaluation protocol.

TranSenseFuser~\cite{transensefuser2023} proposes a hybrid CNN-Transformer
for wrist sensor fusion but likewise evaluates without strict LOSO-CV.

Our work complements these studies by combining a rigorous LOSO-CV protocol,
wrist-only signals, and a systematic ablation that isolates each modality's
contribution across multiple model families---a combination not previously
reported in the literature.

\subsection{EDA in Stress and Affective Computing}

Electrodermal activity reflects sympathetic nervous system arousal via
eccrine sweat gland activation~\cite{boucsein2012}.
EDA has been used extensively for stress detection~\cite{setz2010},
emotion recognition~\cite{koelstra2012}, and cognitive load estimation.
However, EDA responds to \emph{any} sympathetic activation, not stress
specifically: vigorous exercise elicits EDA through thermoregulatory
sweating~\cite{kellogg1994exercise}, and vasomotor hot flashes in
perimenopausal and postmenopausal women produce characteristic
EDA spikes~\cite{freedman2014hotflash}. Context-aware interpretation
of EDA is therefore essential in real-world deployments.

\subsection{Wearable Sensor Modalities}

Photoplethysmography (BVP/PPG) enables heart rate variability (HRV)
estimation from the wrist. HRV is a well-validated autonomic biomarker
of stress~\cite{kim2018hrv}, though wrist PPG quality is lower than
chest ECG for HRV feature extraction.
ACC provides motion context; its primary value in stress detection is
artifact rejection (identifying exercise, ambulatory activity) rather
than direct stress discrimination~\cite{de2019wrist}.
Skin temperature decreases with sympathetic arousal due to peripheral
vasoconstriction, providing a secondary biomarker that is most reliable
under controlled conditions~\cite{vinkers2013temp}.

\section{Dataset and Preprocessing}
\label{sec:dataset}

\subsection{WESAD Dataset}

We use the WESAD (Wearable Stress and Affect Detection) dataset~\cite{schmidt2018wesad},
which contains multimodal physiological recordings from 15 subjects
(S2--S17, S12 missing) during three experimental conditions:
baseline (neutral), TSST (Trier Social Stress Test, cognitive/social stress),
and amusement (Stroop task subset). We adopt binary classification:
TSST windows labeled stress ($y=1$) versus baseline windows ($y=0$),
discarding amusement and transition windows.

Wrist signals are recorded on an Empatica E4:
\begin{itemize}
  \item \textbf{BVP}: 64\,Hz blood volume pulse (PPG)
  \item \textbf{ACC}: 32\,Hz three-axis accelerometry (1/64\,g per count)
  \item \textbf{EDA}: 4\,Hz electrodermal activity ($\mu$S)
  \item \textbf{TEMP}: 4\,Hz skin temperature (\textdegree C)
\end{itemize}

\subsection{Windowing}

We apply non-overlapping sliding windows of 60\,s duration with 30\,s stride,
producing approximately 80--93 windows per subject depending on recording length.
The 60\,s window is a standard choice for HRV feature extraction (capturing
multiple cardiac cycles and low-frequency HRV components) and provides
sufficient temporal context for EDA tonic and phasic components.

ACC is upsampled from 32\,Hz to 64\,Hz (matching BVP) via linear interpolation
for CNN input. EDA and TEMP are upsampled from 4\,Hz to 64\,Hz for CNN input.
For RF feature extraction, all signals are processed at their native sampling rates.

\subsection{Artifact Rejection}

We exclude windows with BVP signal quality below an empirical threshold
(signal power $< 10^{-6}$, corresponding to flat or clipped PPG).
Approximately 2--5\% of windows are discarded per subject.
Remaining windows are class-balanced at the subject level: the stress
class accounts for 23--27\% of valid windows across subjects.

\section{Models}
\label{sec:models}

\subsection{Random Forest with Handcrafted Features}

We extract 34 features per 60\,s window grouped by modality:

\textbf{HRV features (15):} From BVP, we detect R-peaks using a
band-pass filtered (0.5--3\,Hz), normalized, threshold-based peak detector.
Time-domain HRV: mean HR, HR std, SDNN, RMSSD, pNN50.
Poincar\'e plot: SD1, SD2, SD1/SD2. Frequency-domain via Lomb-Scargle
periodogram: LF power (0.04--0.15\,Hz), HF power (0.15--0.4\,Hz),
LF/HF ratio, total power. PPG morphology: mean amplitude, mean rise time,
mean pulse width.

\textbf{ACC features (9):} Magnitude mean and std, signal magnitude area (SMA),
per-axis (x, y, z) mean and std.

\textbf{EDA features (5):} Mean, std, min, max, linear slope.

\textbf{TEMP features (5):} Mean, std, min, max, linear slope.

The RF is a 300-tree ensemble with balanced class weights, trained with
\texttt{sklearn}'s \texttt{RandomForestClassifier}. NaN features
(caused by insufficient peaks in short segments) are imputed with
per-feature training-set means.

\subsection{1D Convolutional Neural Network}

Our CNN follows a three-block architecture:

\textbf{Input:} $(B, C, T)$ where $C \in \{1,2,4,5,6\}$ channels and
$T=3840$ time steps (60\,s $\times$ 64\,Hz).

\textbf{ConvBlocks 1--3:} Each block applies Conv1d $\to$ BatchNorm $\to$ ReLU $\to$
MaxPool(2). Filter counts are $32 \to 64 \to 128$; kernel size 7 with
padding 3; pooling halves the time dimension at each stage ($3840 \to
1920 \to 960 \to 480$).

\textbf{Classifier:} AdaptiveAvgPool1d(1) $\to$ Flatten $\to$ Dropout(0.5) $\to$
Linear(128, 2).

Parameter count: $\approx 140$\,K (independent of $C$).

Training uses AdamW (lr=$10^{-3}$, weight decay=$10^{-4}$), class-weighted
cross-entropy, cosine learning rate annealing, and early stopping on
validation F1 with patience=10. Batch size 32, up to 30 epochs per fold.

\subsection{Patch-based Transformer}

Inspired by PatchTST~\cite{nie2023patchtst} and ViT~\cite{dosovitskiy2020vit},
we design a patch-based Transformer adapted for multimodal 1D physiological signals.

\textbf{Patch Embedding:} The input $(B, C, T)$ is split into non-overlapping
patches of $P=64$ samples (1\,s at 64\,Hz), covering one cardiac cycle.
Each patch flattens all $C$ channels: patch dimension $= C \times P$.
A linear projection maps to $d_\text{model}=64$, followed by LayerNorm,
yielding $n_p = T/P = 60$ patch tokens.

\textbf{CLS Token:} A learnable CLS token is prepended,
giving sequence length $1 + n_p = 61$. Learnable positional embeddings
(truncated normal init, $\sigma=0.02$) are added.

\textbf{Encoder:} Three Pre-LayerNorm TransformerEncoder layers, each with
$n_\text{heads}=4$ attention heads, $d_\text{ff}=128$ feedforward dimension,
and $\text{dropout}=0.3$.

\textbf{Classifier:} The CLS token output passes through Dropout(0.3) $\to$
Linear(64, 2).

Parameter count: $\approx 180$\,K.

Training uses AdamW (lr=$5\times10^{-4}$, weight decay=$10^{-4}$), gradient
clipping (max norm=1.0), class-weighted cross-entropy, cosine annealing,
and early stopping with patience=12. Batch size 32, up to 40 epochs.

The patch size of 1\,s aligns with cardiac cycle length at rest ($\approx$50--100\,bpm),
allowing self-attention to learn which seconds within the 60\,s window
co-vary under stress.

\section{Experimental Setup}
\label{sec:setup}

All experiments use Leave-One-Subject-Out Cross-Validation: for each of
the 15 subjects, that subject is held out as the test set while the
remaining 14 subjects form the training set. For CNN and Transformer,
10\% of training windows are held out as a validation set for early stopping.
No data from the test subject is used in any training or validation step.

The evaluation metric is binary F1 score (stress class), which accounts for
class imbalance (approximately 25\% stress windows). We also report
AUC-ROC and accuracy. Results are summarized as mean $\pm$ std across
the 15 LOSO folds.

All experiments run on CPU. Code and results are available at
\url{https://github.com/ranganarayan/physiological-signal-ml}.

\section{Results}
\label{sec:results}

\subsection{Main Model Comparison}

Table~\ref{tab:main} reports per-subject and aggregate performance for
the three model families, all using the full six-channel input
(BVP, ACC$_{xyz}$, TEMP, EDA).

\begin{table}[t]
  \centering
  \caption{LOSO-CV results on WESAD (15 subjects), full 6-channel input.
           RF: 34 handcrafted features. CNN and Transformer: raw signals.
           F1 = binary stress F1. Best per-row in \textbf{bold}.}
  \label{tab:main}
  \small
  \begin{tabular}{lrrr}
    \toprule
    Subject & RF & CNN & Transformer \\
    \midrule
    S2  & \textbf{1.000} & 0.950 & \textbf{1.000} \\
    S3  & 0.842 & 0.977 & \textbf{0.900} \\
    S4  & 0.973 & 0.944 & \textbf{1.000} \\
    S5  & 0.920 & \textbf{1.000} & \textbf{1.000} \\
    S6  & 0.927 & \textbf{1.000} & 0.974 \\
    S7  & 0.930 & \textbf{1.000} & \textbf{1.000} \\
    S8  & \textbf{1.000} & \textbf{1.000} & \textbf{1.000} \\
    S9  & 0.894 & \textbf{1.000} & \textbf{1.000} \\
    S10 & \textbf{1.000} & \textbf{1.000} & \textbf{1.000} \\
    S11 & 0.914 & 0.744 & \textbf{0.750} \\
    S13 & 0.933 & \textbf{1.000} & \textbf{1.000} \\
    S14 & 0.000 & 0.167 & \textbf{0.326} \\
    S15 & 0.978 & \textbf{1.000} & \textbf{1.000} \\
    S16 & 0.955 & \textbf{1.000} & \textbf{1.000} \\
    S17 & 0.200 & \textbf{0.680} & 0.444 \\
    \midrule
    \textbf{Mean} & 0.831 & \textbf{0.897} & 0.893 \\
    \textbf{Std}  & 0.292 & 0.218 & 0.211 \\
    \textbf{AUC}  & \textbf{0.953} & 0.943 & 0.945 \\
    \bottomrule
  \end{tabular}
\end{table}

The CNN and Transformer achieve nearly identical mean F1 (0.897 and 0.893),
followed by RF (0.831). To test these differences we apply exact Wilcoxon
signed-rank tests on the paired per-subject F1 scores (two-sided, $n=15$).
The CNN--Transformer difference is not significant ($W=14$, $p=1.00$),
whereas both deep models outperform RF (RF vs.\ Transformer $W=10$,
$p=0.021$; RF vs.\ CNN $W=18$, $p=0.057$). Subject-level bootstrap 95\%
confidence intervals for the mean F1 are RF $[0.662, 0.952]$,
CNN $[0.770, 0.987]$, and Transformer $[0.774, 0.985]$; the wide intervals
reflect the small 15-subject sample.
Importantly, the RF, despite relying on handcrafted features, performs
competitively (only 6.2 F1 points below the Transformer), indicating
that a well-tuned feature engineering pipeline captures most of the
discriminative information available in 60\,s windows.

Subjects S14 and S17 are consistent outliers across all three models
(Section~\ref{sec:subjectvar}).

\subsection{CNN Channel Ablation}
\label{sec:cnn_ablation}

Table~\ref{tab:cnn_ablation} reports CNN LOSO-CV F1 across six channel
configurations. We systematically add modalities to a BVP baseline.

\begin{table}[t]
  \centering
  \caption{CNN channel ablation. Mean F1 $\pm$ std across 15 LOSO folds.
           The ablation is an independent set of LOSO trainings; its full
           six-channel entry (0.878) differs from the main-comparison CNN
           in Table~\ref{tab:main} (0.897) by run-to-run stochastic variance,
           itself indicative of the training noise at $n=15$.}
  \label{tab:cnn_ablation}
  \small
  \setlength{\tabcolsep}{4pt}
  \resizebox{\columnwidth}{!}{%
  \begin{tabular}{lrrr}
    \toprule
    Configuration & Channels & Mean F1 & Std \\
    \midrule
    BVP only          & 1 & 0.737 & 0.094 \\
    BVP + ACC         & 4 & 0.589 & 0.175 \\
    BVP + ACC + TEMP  & 5 & 0.599 & 0.212 \\
    BVP + ACC + EDA   & 5 & 0.884 & 0.154 \\
    BVP + EDA         & 2 & 0.857 & 0.179 \\
    BVP + ACC + TEMP + EDA & 6 & 0.878 & 0.232 \\
    \bottomrule
  \end{tabular}}
\end{table}

Several findings stand out:

\textbf{EDA is the dominant channel.} Adding EDA to the BVP+ACC baseline
yields $+0.295$ F1 (0.589 $\to$ 0.884), while adding TEMP yields only $+0.010$
(0.589 $\to$ 0.599). The marginal value of EDA is approximately 30$\times$
larger than that of TEMP.

\textbf{ACC hurts in the sedentary TSST context.} BVP-only (0.737) outperforms
BVP+ACC (0.589), a $-0.148$ F1 degradation. The WESAD TSST protocol is
conducted seated; ACC signals are near-constant (low-magnitude noise around
baseline). Adding ACC thus contributes only noise with no discriminative content.

\textbf{Two channels suffice.} BVP+EDA without ACC or TEMP achieves 0.857,
within 2.1 F1 points of the full six-channel model (0.878). This finding
has direct implications for wearable sensor design: a device with only an
optical heart rate sensor and an EDA electrode can capture nearly all
available stress-discriminative information in sedentary conditions.

\textbf{TEMP adds noise when EDA is present.} BVP+ACC+EDA (0.884) vs.
BVP+ACC+TEMP+EDA (0.878) shows TEMP slightly decreases mean F1 when
all other modalities are included, largely driven by a catastrophic
interaction in subject S14 (see Section~\ref{sec:subjectvar}).

\subsection{RF Feature Group Ablation}
\label{sec:rf_ablation}

Table~\ref{tab:rf_ablation} reports RF performance for ten handcrafted
feature subsets. This analysis tests whether EDA dominance reflects the
physiological signal itself or is an artifact of CNN architecture.

\begin{table}[t]
  \centering
  \caption{RF feature group ablation. Mean F1 $\pm$ std across 15 LOSO folds.}
  \label{tab:rf_ablation}
  \small
  \setlength{\tabcolsep}{4pt}
  \resizebox{\columnwidth}{!}{%
  \begin{tabular}{lrrr}
    \toprule
    Feature Group & Features & Mean F1 & Std \\
    \midrule
    EDA only         &  5 & 0.774 & 0.234 \\
    HRV only         & 15 & 0.622 & 0.118 \\
    TEMP only        &  5 & 0.489 & 0.227 \\
    ACC only         &  9 & 0.323 & 0.140 \\
    \midrule
    EDA + TEMP       & 10 & 0.804 & 0.316 \\
    HRV + EDA        & 20 & 0.829 & 0.226 \\
    HRV + TEMP       & 20 & 0.628 & 0.130 \\
    HRV + ACC        & 24 & 0.669 & 0.123 \\
    \midrule
    HRV + EDA + TEMP & 25 & 0.826 & 0.296 \\
    All (34 features)& 34 & \textbf{0.831} & 0.292 \\
    \bottomrule
  \end{tabular}}
\end{table}

The RF ablation corroborates and extends the CNN findings:

\textbf{EDA alone (0.774) outperforms HRV alone (0.622)} by 15.2 F1 points.
Five simple statistics (mean, std, min, max, slope) over the EDA signal
carry more stress-discriminative information than 15 carefully crafted HRV
features including time-domain, frequency-domain, and Poincar\'e metrics.

\textbf{ACC is the weakest modality} in isolation (0.323), consistent with
the CNN ablation result.

\textbf{HRV + EDA (0.829) nearly matches All features (0.831)}, with
$\Delta = 0.002$ for adding nine ACC and five TEMP features.
This confirms that the informative content of the 34-feature set is
concentrated in the 20 HRV and EDA features.

\textbf{EDA dominance is architecture-independent.}
The finding appears in both the end-to-end CNN (raw signals) and the
RF (handcrafted features), ruling out a CNN-specific modeling artifact.

\begin{figure*}[t]
  \centering
  \includegraphics[width=\textwidth]{ablation_figure.pdf}
  \caption{Channel and feature ablation results under LOSO-CV (15 subjects).
           Error bars show $\pm$1 std across folds.
           \textbf{(a) CNN channel ablation:} EDA is the dominant modality;
           BVP+EDA (2-channel) achieves 0.857 F1, within 2.1 points of the
           full 6-channel model. ACC degrades performance in the sedentary TSST context.
           \textbf{(b) RF feature group ablation:} EDA alone (0.774) outperforms
           HRV alone (0.622) by 15.2 F1 points, confirming EDA dominance is
           architecture-independent. Green bars indicate EDA-containing configurations.}
  \label{fig:ablation}
\end{figure*}

\subsection{Per-Subject Variability}
\label{sec:subjectvar}

Table~\ref{tab:subject_ablation} reports per-subject CNN F1 across five
key channel configurations, making the heterogeneous impact of each
modality across individuals explicit.

\begin{table}[t]
  \centering
  \caption{Per-subject CNN F1 for five channel configurations.
           B+E = BVP+EDA (2ch); B+A+E = BVP+ACC+EDA; 6ch = all channels.
           Mean row reproduces Table~\ref{tab:cnn_ablation}.}
  \label{tab:subject_ablation}
  \footnotesize
  \setlength{\tabcolsep}{4pt}
  \begin{tabular}{lrrrrr}
    \toprule
    S & BVP & B+A & B+E & B+A+E & 6ch \\
    \midrule
    S2  & 0.706 & 0.462 & 0.884 & 0.864 & 0.950 \\
    S3  & 0.731 & 0.525 & 0.704 & 0.840 & 0.894 \\
    S4  & 0.737 & 0.818 & 0.944 & 0.944 & 1.000 \\
    S5  & 0.755 & 0.615 & 0.852 & 0.902 & 1.000 \\
    S6  & 0.818 & 0.882 & 0.864 & 0.905 & 0.974 \\
    S7  & 0.766 & 0.708 & 0.930 & 0.952 & 0.909 \\
    S8  & 0.809 & 0.232 & 1.000 & 0.952 & 1.000 \\
    S9  & 0.582 & 0.737 & 0.977 & 1.000 & 1.000 \\
    S10 & 0.750 & 0.558 & 1.000 & 1.000 & 1.000 \\
    S11 & 0.760 & 0.609 & 0.727 & 0.783 & 0.762 \\
    S13 & 0.833 & 0.630 & 0.913 & 1.000 & 1.000 \\
    S14 & 0.727 & 0.350 & 0.769 & 0.732 & 0.138 \\
    S15 & 0.484 & 0.353 & 1.000 & 1.000 & 1.000 \\
    S16 & 0.880 & 0.656 & 1.000 & 1.000 & 1.000 \\
    S17 & 0.720 & 0.704 & 0.290 & 0.393 & 0.549 \\
    \midrule
    Mean & 0.737 & 0.589 & 0.857 & 0.884 & 0.878 \\
    \bottomrule
  \end{tabular}
\end{table}

\textbf{Subject S14.} S14 is the worst-performing subject across all
three models (CNN: 0.167, RF: 0.000, Transformer: 0.326,
Table~\ref{tab:main}). The per-channel ablation reveals a striking
pattern: S14 achieves 0.727 with BVP alone, 0.769 with BVP+EDA,
but collapses to 0.138 under the full six-channel input.
Intermediate configurations confirm that TEMP is the destructive modality:
adding TEMP alongside EDA drives F1 toward zero. This suggests S14 has
an atypical thermoregulatory response during the TSST such that their
skin temperature signature is inversely correlated with the stress-class
TEMP pattern learned from the other 14 subjects. RF similarly fails on
S14 (F1=0.000), confirming the issue is in the physiological signal
rather than a modeling artifact.

\textbf{Subject S17.} S17 exhibits an \emph{inverted EDA response}.
BVP alone yields F1=0.720; adding EDA (BVP+EDA) drops performance to
0.290, a $-0.430$ F1 degradation from a single modality addition.
Adding ACC alongside EDA (BVP+ACC+EDA) recovers partially to 0.393,
and the full six-channel model reaches 0.549, still well below the
BVP-only baseline. The pattern is consistent across all three model
families: RF achieves 0.200 and Transformer 0.444 on S17, both
substantially below their respective population means.
This is consistent with a suppressed or inverted sympathetic EDA response
to psychological stressors, a documented phenomenon in a small fraction
of the population, potentially related to habituation, atypical autonomic
regulation, or EDA sensor contact quality.

\textbf{Subject S15.} S15 illustrates the opposite pattern: BVP alone
yields a below-average 0.484 F1, but adding EDA rescues performance
completely (BVP+EDA = 1.000, BVP+ACC+EDA = 1.000, 6ch = 1.000).
For S15, HRV-derived features alone are insufficient, but EDA provides
a clean stress signal that fully disambiguates stress from baseline.

\textbf{S8 ACC degradation.} S8's BVP+ACC F1 drops to 0.232 from a
BVP-only baseline of 0.809, the most severe ACC-induced degradation
in the dataset. Adding EDA restores S8 to 1.000 (BVP+EDA and 6ch).
This suggests S8 had unusual motion patterns during the TSST that the
model mistakenly weighted as non-stress signal.

\textbf{High-performing subjects.} Subjects S5--S10, S13, S15, and S16
reach F1 $\geq 0.95$ in the BVP+EDA or full six-channel configuration,
confirming the task is highly learnable for the majority of the population.

\section{Discussion}
\label{sec:discussion}

\subsection{EDA as a Dominant but Context-Sensitive Biomarker}

EDA is the single most informative wrist-sensor modality for stress detection
in the WESAD TSST protocol, which reliably elicits sympathetic
activation---public speaking and mental arithmetic before an evaluating
panel---that manifests as eccrine sweat-gland activity at the wrist.

However, EDA reflects \emph{sympathetic arousal broadly}, not stress
specifically. Three confounds matter for real-world deployment:

\textbf{Exercise:} Vigorous activity drives thermoregulatory sweating that is
physiologically distinct from stress-related palmar sweating but electrically
indistinguishable at the sensor~\cite{kellogg1994exercise}. We propose
ACC-based gating: when ACC root-mean-square exceeds an activity threshold,
stress classification is suppressed. WESAD subjects are sedentary, but a
deployed system would apply this gate in real time.

\textbf{Thermal environment:} Ambient temperatures above $\sim$28\textdegree C
elicit sweating even at rest, confounding EDA~\cite{romanovsky2014thermoreg};
a wrist skin-temperature channel can flag this as a context input.

\textbf{Vasomotor events:} Hot flashes produce characteristic EDA spikes via
hypothalamic thermoregulatory dysregulation~\cite{freedman2014hotflash}, so
stress monitoring in perimenopausal and postmenopausal populations would
require demographically stratified models or explicit vasomotor context.

\subsection{ACC as Context, Not Signal}

Accelerometry's negative contribution in WESAD reflects the seated TSST
protocol. In free-living use, ACC is valuable as a \emph{context signal}
(separating exercise from rest) rather than a direct stress biomarker, and
future work should test it on datasets with mixed activity profiles.
Disambiguating its role under simultaneous physical and cognitive load---relevant
for athletes and emergency responders---would require dual-task paradigms such
as mental arithmetic during exercise.

\subsection{Implications for Recovery and Allostatic Load Monitoring}

Our finding has implications beyond acute stress alerting. Recovery-focused
wearables model physiological load from HR and HRV, using overnight HRV
suppression as a proxy for incomplete recovery---capturing \emph{physical}
strain well. But cognitive stress activates the same
hypothalamic-pituitary-adrenal axis and sympathetic pathways~\cite{kim2018hrv},
producing autonomic load invisible to activity-based metrics: a sedentary
knowledge worker under sustained stress may show low physical strain yet
suppressed overnight HRV that a load model then misattributes to noise.

EDA offers a direct window into that sympathetic arousal precisely where
ACC-based strain is near-zero. Since even five tonic EDA statistics outperform
a 15-feature HRV set, intermittent sampling (e.g., 60\,s every few minutes)
could estimate cumulative psychological load across the waking day at modest
power cost---adding a psychological component to physical-strain-based recovery
estimates without continuous high-frequency streaming.

\subsection{Architecture Selection}

Our results show that all three architectures perform comparably on WESAD,
with the CNN and Transformer statistically tied at the top (0.897 and 0.893
mean F1; Wilcoxon $p=1.00$) and RF close behind (0.831). The practical implication depends on deployment context:

\textbf{Edge inference:} RF (34 features, tree evaluation) runs in microseconds
on an embedded processor. The CNN ($\approx$140\,K parameters) requires a
small neural accelerator. The Transformer ($\approx$180\,K parameters) is
slightly heavier but still feasible on modern smartwatch SoCs.

\textbf{Training data efficiency:} With 15 subjects and $\approx$1,300
training windows per fold, the RF and CNN learn stably. The Transformer
converges more slowly (patience=12 vs.\ 10 for CNN) and may benefit from
larger datasets. We expect the Transformer advantage to grow with $N>100$ subjects.

\textbf{Label efficiency:} The Transformer's attention mechanism allows it
to learn which seconds within the 60\,s window co-vary under stress,
potentially capturing temporal patterns (stress onset, sustained elevation)
that the CNN's fixed-receptive-field architecture misses.

\subsection{Scalability and Personalization}

All models are trained from scratch per LOSO fold. How performance scales
with cohort size is open; we expect the Transformer to gain most as $N$ grows
toward hundreds or thousands, given the data appetite of attention models.
For outliers like S14 and S17, subject-specific fine-tuning is promising---even
5--10 labeled windows could supply a personalization signal without full
per-subject training.

\subsection{Limitations}

\textbf{Dataset size:} 15 subjects is small for cross-subject generalization;
mean LOSO F1 estimates have high variance ($\text{std} \approx 0.21$--$0.29$).

\textbf{Lab protocol:} The TSST is a controlled acute stressor; chronic stress,
occupational stress, and naturalistic daily stress have different physiological
signatures.

\textbf{Population:} WESAD subjects are healthy young adults. Generalization
to older populations, individuals with cardiovascular conditions, or those
with unusual autonomic profiles (as seen in S14 and S17) is uncertain.

\textbf{EDA power consumption:} Continuous EDA sensing requires a small
constant current through the skin, adding power draw relative to PPG-only
devices. This is a practical constraint for wearables targeting multi-day
battery life. Intermittent EDA sampling (60\,s epochs every 5--10 minutes)
is a feasible mitigation, as our tonic-level features (mean, std, slope)
do not require continuous high-frequency streaming and sustained
stress episodes persist across multiple sampling windows.

\textbf{Limited statistical power:} We report exact Wilcoxon signed-rank
tests and subject-level bootstrap confidence intervals
(Section~\ref{sec:results}), but with $n=15$ paired folds these have
limited power to detect small effects. The non-significant
CNN--Transformer comparison should therefore be read as ``no detectable
difference'' rather than proven equivalence.

\section{Conclusion}
\label{sec:conclusion}

We conducted a systematic comparison of Random Forest, 1D CNN, and patch-based
Transformer architectures for physiological stress detection on WESAD under
strict Leave-One-Subject-Out Cross-Validation. All three models use identical
six-channel wrist-sensor input (BVP, ACC, TEMP, EDA) and identical data splits,
enabling a fair architectural comparison. The CNN and Transformer perform
comparably (0.897 and 0.893 mean F1; Wilcoxon signed-rank $p=1.00$), both
ahead of RF (0.831); the CNN--Transformer difference is not statistically
significant at $n=15$.

Our channel ablation study is the paper's primary contribution. EDA is the
dominant stress biomarker across both CNN and RF architectures: adding EDA
to a BVP$+$ACC baseline improves CNN F1 by $+0.295$ ($0.589 \to 0.884$),
and in the RF, EDA-only features (0.774) outperform HRV-only (0.622).
Accelerometry \emph{hurts}
performance in sedentary conditions ($-0.148$ F1), and skin temperature
provides negligible marginal information once EDA is included. A two-channel
BVP+EDA configuration achieves 0.857 F1, only 2.1 points below the full
six-channel model, with practical implications for wearable sensor design:
wrist devices need not include ACC or TEMP to achieve near-peak stress
detection performance in sedentary contexts.

Per-subject analysis reveals clinically important heterogeneity: two subjects
(S14, S17) are consistent outliers attributable to atypical thermoregulatory
(S14) and electrodermal (S17) responses, motivating subject-aware and
personalized stress models.

Future directions include: (1) evaluation on larger and more diverse datasets,
(2) ACC-gated deployment for active populations, (3) subject-adaptive
fine-tuning for outlier subjects, and (4) extension to chronic and
naturalistic stress paradigms.

\section*{Acknowledgments}
The author thanks the WESAD dataset creators at the University of Siegen
for making the benchmark publicly available.

\bibliographystyle{plain}
\begin{thebibliography}{99}

\bibitem{schmidt2018wesad}
P.~Schmidt, A.~Reiss, R.~Duerichen, C.~Marberger, and K.~Van~Laerhoven,
``Introducing WESAD, a multimodal dataset for wearable stress and affect
detection,'' in \textit{Proc. ACM ICMI}, 2018.

\bibitem{boucsein2012}
W.~Boucsein, \textit{Electrodermal Activity}, 2nd ed. Springer, 2012.

\bibitem{setz2010}
C.~Setz, B.~Arnrich, J.~Schumm, R.~La~Marca, G.~Tr\"oster, and U.~Ehlert,
``Discriminating stress from cognitive load using a wearable EDA device,''
\textit{IEEE Trans. Inf. Technol. Biomed.}, vol.~14, no.~2, pp.~410--417, 2010.

\bibitem{koelstra2012}
S.~Koelstra et al., ``DEAP: A database for emotion analysis using physiological
signals,'' \textit{IEEE Trans. Affect. Comput.}, vol.~3, no.~1, pp.~18--31, 2012.

\bibitem{kim2018hrv}
H.-G.~Kim, E.-J.~Cheon, D.-S.~Bai, Y.~H.~Lee, and B.-H.~Koo,
``Stress and heart rate variability: A meta-analysis and review of the
literature,'' \textit{Psychiatry Investig.}, vol.~15, no.~3, pp.~235--245, 2018.

\bibitem{yaribeygi2017}
H.~Yaribeygi, Y.~Panahi, H.~Sahraei, T.~P.~Johnston, and A.~Sahebkar,
``The impact of stress on body function: A review,''
\textit{EXCLI J.}, vol.~16, pp.~1057--1072, 2017.

\bibitem{mekruksavanich2021}
S.~Mekruksavanich and A.~Jitpattanakul,
``LSTM networks for physiological stress detection,''
\textit{IEEE Access}, 2021.

\bibitem{cross_modality_2025}
E.~Oliver and S.~Dakshit,
``Cross-modality investigation on WESAD stress classification,''
arXiv preprint arXiv:2502.18733, 2025.

\bibitem{transensefuser2023}
P.~Kasnesis, C.~Chatzigeorgiou et al.,
``TranSenseFusers: A temporal CNN-Transformer neural network family for
explainable PPG-based stress detection,''
\textit{Biomed. Signal Process. Control}, 2024.

\bibitem{de2019wrist}
M.~de~Zambotti et al.,
``Wrist-worn wearables for measuring heart rate and heart rate variability,''
\textit{NPJ Digit. Med.}, 2019.

\bibitem{vinkers2013temp}
C.~H.~Vinkers et al.,
``Distal-to-proximal skin temperature gradient as a measure of stress,''
\textit{Psychophysiology}, 2013.

\bibitem{kellogg1994exercise}
D.~L.~Kellogg, ``In vivo mechanisms of cutaneous vasodilation
and vasoconstriction in humans during thermoregulatory challenges,''
\textit{J. Appl. Physiol.}, vol.~100, pp.~1709--1718, 2006.

\bibitem{freedman2014hotflash}
R.~R.~Freedman,
``Menopausal hot flashes: mechanisms, endocrinology, treatment,''
\textit{J. Steroid Biochem. Mol. Biol.}, vol.~142, pp.~115--120, 2014.

\bibitem{romanovsky2014thermoreg}
A.~A.~Romanovsky,
``Skin temperature: its role in thermoregulation,''
\textit{Acta Physiol.}, vol.~210, no.~3, pp.~498--507, 2014.

\bibitem{dosovitskiy2020vit}
A.~Dosovitskiy et al.,
``An image is worth 16$\times$16 words: Transformers for image recognition
at scale,'' in \textit{Proc. ICLR}, 2021.

\bibitem{nie2023patchtst}
Y.~Nie, N.~H.~Nguyen, P.~Sinthong, and J.~Kalagnanam,
``A time series is worth 64 words: Long-term forecasting with transformers,''
in \textit{Proc. ICLR}, 2023.

\end{thebibliography}

\end{document}
